% Part: normal-modal-logic
% Chapter: completeness
% Section: modalities-ccs
%READERNOTE{SOL6-A165: Canonical accessibility चे iff आणि शून्य गृहीतकांचा boxing lemma. Diamond branch मधील explanatory clause ने R खरे होण्यासाठी sufficient अट दिली होती; अपेक्षित implication साठी R ते condition ही दिशा आवश्यक आहे. औपचारिक पुढील definition प्रमाणे iff स्पष्ट केला. Boxing lemma च्या k=0 प्रकरणात RK नव्हे तर NEC थेट लागू केला; k>=1 साठी मूळ RK proof जपला. Modal equivalence च्या वाक्यात missing शी जोडले. पूर्वीचे k/n आणि system-subscript repairs योग्य.}

\documentclass[../../../include/open-logic-section]{subfiles}

\begin{document}

\olfileid{nml}{com}{mod}

\olsection{मोडल कारक आणि संपूर्ण सुसंगत संच}

\begin{explain}
एखाद्या सामान्य मोडल तर्कशास्त्र~$\Sigma$ मधील संपूर्ण
$\Sigma$-सुसंगत संच~$\Delta$ हे ज्याचे जग आहेत,
असे प्रतिमान~$\mModel{M^\Sigma}$ रचताना अशा “जगां”मधील
प्राप्यता संबंध~$R^\Sigma$ देखील परिभाषित करावा लागेल.
प्राप्यता संबंध (आणि मूल्यांकन~$V^\Sigma$) यांची व्याख्या
अशी हवी की $\mSat{M^\Sigma}{!A}[\Delta]$ असेल तेव्हा
आणि तेव्हाच $!A \in \Delta$. हे कसे करावे?

प्राप्यता संबंध परिभाषित केल्यावर, जगातील सत्यतेच्या
व्याख्येनुसार \iftag{prvBox}{$\mSat{M^\Sigma}{\Box!A}[\Delta]$
  असेल तेव्हा आणि तेव्हाच $R^\Sigma\Delta\Delta'$ असलेल्या
  प्रत्येक $\Delta'$ साठी $\mSat{M^\Sigma}{!A}[\Delta']$}
  {$\mSat{M^\Sigma}{\Diamond!A}[\Delta]$ असेल तेव्हा
  आणि तेव्हाच $R^\Sigma\Delta\Delta'$ असलेल्या किमान एका
  $\Delta'$ साठी $\mSat{M^\Sigma}{!A}[\Delta']$}.
$\mSat{M^\Sigma}{!A}[\Delta]$ असेल तेव्हा आणि तेव्हाच
$!A \in \Delta$ हे सिद्ध करण्यासाठी हे विशेषतः
मोडल कारकाने सुरू होणाऱ्या !!{formula}s साठी खरे असणे
आवश्यक आहे; म्हणजे
\iftag{prvBox}{$\mSat{M^\Sigma}{\Box!A}[\Delta]$ असेल
  तेव्हा आणि तेव्हाच $\Box !A \in \Delta$}
  {$\mSat{M^\Sigma}{\Diamond!A}[\Delta]$ असेल
  तेव्हा आणि तेव्हाच $\Diamond !A \in \Delta$}.
ही आवश्यकता आणि
\iftag{prvBox}{$\Box !A$}{$\Diamond !A$}
साठी जगातील सत्यतेची व्याख्या एकत्र केल्यास:
\[
\iftag{prvBox}{%
  \Box!A \in \Delta \text{ तेव्हा आणि तेव्हाच } !A \in \Delta' \text{
    प्रत्येक $\Delta'$ साठी, जेथे } R^\Sigma\Delta\Delta'}{%
  \Diamond!A \in \Delta \text{ तेव्हा आणि तेव्हाच } !A \in \Delta' \text{
    किमान एका } \Delta' \text{ साठी, जेथे } R^\Sigma\Delta\Delta'}
\]
\iftag{prvBox}{डावीकडून उजवीकडे जाणारी दिशा पाहा:
  $\Box !A \in \Delta$ असेल, तर कोणत्याही $!A$
  आणि $R^\Sigma\Delta\Delta'$ असलेल्या कोणत्याही
  $\Delta'$ साठी $!A \in \Delta'$ असे ती सांगते.
  प्रत्येक $\Box!A \in \Delta$ साठी $!A \in \Delta'$
  असेल तेव्हा आणि तेव्हाच $R^\Sigma\Delta\Delta'$
  अशी अट ठेवल्यास हे खरे ठरते. या द्विसशर्ताच्या
  उजव्या बाजूची अट संक्षिप्तपणे
  $\Setabs{!A}{\Box!A \in \Delta} \subseteq \Delta'$
  अशी लिहिता येते.}{उजवीकडून डावीकडे जाणारी दिशा
  पाहा: $R^\Sigma\Delta\Delta'$ असलेल्या कोणत्याही
  $!A$ आणि $\Delta'$ साठी $!A \in \Delta'$ असेल,
  तर $\Diamond!A \in \Delta$ असे ती सांगते.
  प्रत्येक $!A \in \Delta'$ साठी $\Diamond!A \in
  \Delta$ असेल तेव्हा आणि तेव्हाच $R^\Sigma\Delta\Delta'$
  खरा ठरतो अशी अट ठेवल्यास हे खरे ठरते. या
  द्विसशर्ताच्या उजव्या बाजूची अट संक्षिप्तपणे
  $\Setabs{\Diamond!A}{!A \in \Delta'} \subseteq \Delta$
  अशी लिहिता येते.}

मग प्रश्न असा आहे: $R^\Sigma$ ची ही व्याख्या
\iftag{prvBox}{$\Box !A \in \Delta$ असेल तेव्हा आणि
  तेव्हाच $\mSat{M^\Sigma}{\Box !A}[\Delta]$ याची
  खरोखर हमी देते का? \iftag{prvDiamond}{ती
    }{}}{}\iftag{prvDiamond}{$\Diamond !A \in \Delta$
  असेल तेव्हा आणि तेव्हाच
  $\mSat{M^\Sigma}{\Diamond !A}[\Delta]$ याचीही
  हमी देते का?}{} पुढील काही निष्कर्ष हे सिद्ध करतील.
\end{explain}

\begin{defn}
  $\Gamma$ हा !!{formula}s चा संच असेल, तर
  \begin{align*}
    \Box\Gamma & = \Setabs{\Box !B}{!B \in \Gamma}\\
    \Diamond\Gamma & = \Setabs{\Diamond !B}{!B \in \Gamma}\\
    \intertext{आणि}
    \Box^{-1}\Gamma & = \Setabs{!B}{\Box !B \in \Gamma}\\
    \Diamond^{-1}\Gamma & = \Setabs{!B}{\Diamond !B \in \Gamma}\\
  \end{align*}
\end{defn}

दुसऱ्या शब्दांत, $\Box\Gamma$ म्हणजे~$\Gamma$ मधील
प्रत्येक !!{formula} च्या पुढे $\Box$ लावलेला संच;
ही सर्व सूत्रे~$\Gamma$ मधूनच घेतली जातात.
$\Box^{-1}\Gamma$ मध्ये~$\Gamma$ मधील $\Box$ ने
सुरू होणाऱ्या सर्व !!{formula}s मधील सुरुवातीचा
$\Box$ काढून टाकला जातो. ही व्याख्या स्वतंत्रपणे
फारशी महत्त्वाची नाही, पण संकेतलेखन बरेच सोपे करते.

लक्षात घ्या की $\Box\Box^{-1}\Gamma \subseteq \Gamma$:
\[
\Box\Box^{-1}\Gamma = \Setabs{\Box!B}{\Box!B \in \Gamma}
\]
म्हणजे हा~$\Gamma$ मधील $\Box$ ने सुरू होणाऱ्या
सर्व !!{formula}s चाच संच आहे.

\begin{lem}\ollabel{lem:box1}
  $\Gamma \Proves[\Sigma] !A$ असेल, तर
  $\Box\Gamma \Proves[\Sigma] \Box !A$.
\end{lem}

\begin{proof}
  $\Gamma \Proves[\Sigma] !A$ असेल, तर
  $!B_1$, \dots, $!B_k \in \Gamma$ अशी सूत्रे
  असतात. $k=0$ असेल, तर $\Sigma \Proves !A$;
  \Nec{} ने $\Sigma \Proves \Box!A$ मिळते आणि
  शून्य गृहीतकांच्या प्रकरणानुसार $\Box\Gamma \Proves[\Sigma] \Box!A$.
  आता $k\geq1$ असे समजा. मग $\Sigma \Proves !B_1 \lif (!B_2 \lif \cdots
  (!B_k \lif !A)\cdots)$. $\Sigma$ सामान्य असल्याने
  \RK{} नियमानुसार $\Sigma \Proves \Box!B_1 \lif
  (\Box!B_2 \lif \cdots (\Box!B_k \lif \Box!A)\cdots)$;
  येथे अर्थातच $\Box!B_1$, \dots, $\Box!B_k \in
  \Box\Gamma$. म्हणून व्याख्येनुसार
  $\Box\Gamma \Proves[\Sigma] \Box !A$.
\end{proof}

\begin{lem}\ollabel{lem:box2}
  $\Box^{-1} \Gamma \Proves[\Sigma] !A$ असेल,
  तर $\Gamma \Proves[\Sigma] \Box!A$.
\end{lem}

\begin{proof}
  $\Box^{-1}\Gamma \Proves[\Sigma] !A$ असे
  समजा. मग \olref{lem:box1} नुसार
  $\Box\Box^{-1}\Gamma \Proves[\Sigma] \Box!A$.
  पण $\Box\Box^{-1}\Gamma \subseteq \Gamma$
  असल्याने एकस्वनिकतेनुसार
  $\Gamma \Proves[\Sigma] \Box!A$.
\end{proof}

\iftag{prvBox}{%
% Box is primitive

\begin{prop}\ollabel{prop:box}
  $\Gamma$ हा संपूर्ण $\Sigma$-सुसंगत असेल,
  तर $\Box !A \in \Gamma$ असेल तेव्हा आणि
  तेव्हाच $\Box^{-1} \Gamma \subseteq \Delta$
  असलेल्या प्रत्येक संपूर्ण $\Sigma$-सुसंगत
  $\Delta$ साठी $!A \in \Delta$ असेल.
\end{prop}

\begin{proof}
  $\Gamma$ हा संपूर्ण $\Sigma$-सुसंगत आहे
  असे समजा. डावीकडून उजवीकडे जाणारी दिशा सोपी आहे:
  $\Box !A \in \Gamma$ आणि
  $\Box^{-1}\Gamma \subseteq \Delta$ असे
  समजा. $\Box!A \in \Gamma$ असल्याने
  $!A \in \Box^{-1}\Gamma \subseteq \Delta$;
  म्हणून $!A \in \Delta$.

  उजवीकडून डावीकडे जाणाऱ्या दिशेसाठी आपण प्रतिपरिवर्तन सिद्ध करू:
  $\Box!A \notin \Gamma$ असे समजा.
  $\Gamma$ हा संपूर्ण $\Sigma$-सुसंगत
  असल्याने तो निगामीदृष्ट्या बंद आहे;
  म्हणून $\Gamma \Proves/[\Sigma] \Box !A$.
  \olref{lem:box2} नुसार
  $\Box^{-1}\Gamma \Proves/[\Sigma] !A$.
  \olref[prf][con]{prop:consistencyfacts}\olref[prf][con]{prop:consistencyfacts-b}
  नुसार $\Box^{-1}\Gamma \cup \{ \lnot!A \}$
  हा $\Sigma$-सुसंगत आहे. लिंडेनबाउमच्या
  पूर्वप्रमेयानुसार $\Box^{-1}\Gamma \cup
  \{ \lnot!A \} \subseteq \Delta$ असा
  संपूर्ण $\Sigma$-सुसंगत संच~$\Delta$
  असतो. सुसंगततेमुळे $!A \notin \Delta$.
\end{proof}
}{%
% Box is not primitive, only Diamond is

\begin{prop}\ollabel{prop:diamond}
  $\Gamma$ हा संपूर्ण $\Sigma$-सुसंगत असेल,
  तर $\Diamond !A \in \Gamma$ असेल तेव्हा
  आणि तेव्हाच $\Diamond\Delta \subseteq
  \Gamma$ असलेला एखादा संपूर्ण
  $\Sigma$-सुसंगत $\Delta$ असा असेल की
  $!A \in \Delta$.
\end{prop}

\begin{proof}
  $\Gamma$ हा संपूर्ण $\Sigma$-सुसंगत आहे
  असे समजा. उजवीकडून डावीकडे जाणारी दिशा
  सोपी आहे: $\Diamond\Delta \subseteq
  \Gamma$ आणि $!A \in \Delta$ असलेला
  संपूर्ण $\Sigma$-सुसंगत $\Delta$
  घ्या. मग $\Diamond!A \in \Diamond\Delta
  \subseteq \Gamma$; म्हणजे
  $\Diamond!A \in \Gamma$.

  डावीकडून उजवीकडे जाणाऱ्या दिशेसाठी
  $\Diamond !A \in \Gamma$ असे समजा.
  $!A \in \Delta$ आणि
  $\Diamond\Delta \subseteq \Gamma$
  असलेला संपूर्ण $\Sigma$-सुसंगत
  $\Delta$ अस्तित्वात आहे हे दाखवायचे
  आहे. $\Diamond !A \in \Gamma$
  आणि $\Gamma$ $\Sigma$-सुसंगत असल्याने
  $\Box\lnot !A \notin \Gamma$.
  $\Gamma$ हा संपूर्ण $\Sigma$-सुसंगत
  असल्याने तो निगामीदृष्ट्या बंद आहे;
  म्हणून $\Gamma \Proves/[\Sigma]
  \Box\lnot !A$. \olref{lem:box2} नुसार
  $\Box^{-1}\Gamma \Proves/[\Sigma] \lnot !A$.
  \olref[prf][con]{prop:consistencyfacts}\olref[prf][con]{prop:consistencyfacts-b}
  नुसार $\Box^{-1}\Gamma \cup \{ !A \}$
  हा $\Sigma$-सुसंगत आहे. लिंडेनबाउमच्या
  पूर्वप्रमेयानुसार $\Box^{-1}\Gamma \cup
  \{ !A \} \subseteq \Delta$ असा
  संपूर्ण $\Sigma$-सुसंगत संच~$\Delta$
  असतो. स्पष्टच, $!A \in \Delta$.
  $\Diamond\Delta \subseteq \Gamma$
  हे पाहण्यासाठी तसे नाही असे समजा:
  एखाद्या $!B \in \Delta$ साठी
  $\Diamond !B \notin \Gamma$. मग
  $\Gamma$ संपूर्ण $\Sigma$-सुसंगत
  असल्यामुळे $\Box\lnot !B \in \Gamma$.
  पण मग $\lnot !B \in \Box^{-1}\Gamma
  \subseteq \Delta$ हेही मिळते; त्यामुळे
  $\Delta$ विसंगत ठरेल.
\end{proof}
}

\begin{lem}\ollabel{lem:box-iff-diamond}
  $\Gamma$ आणि $\Delta$ हे संपूर्ण
  $\Sigma$-सुसंगत असोत. मग
  $\Box^{-1}\Gamma \subseteq \Delta$
  असेल तेव्हा आणि तेव्हाच
  $\Diamond\Delta \subseteq \Gamma$.
\end{lem}

\begin{proof}
  डावीकडून उजवीकडे जाणारी दिशा: $\Box^{-1}\Gamma
  \subseteq \Delta$ असे समजा आणि
  $\Diamond!A \in \Diamond\Delta$
  (म्हणजे $!A \in \Delta$) असे धरा.
  $\Diamond!A \in \Gamma$ दाखवण्यासाठी
  $\Box\lnot!A \notin \Gamma$ दाखवणे
  पुरेसे आहे; कारण मग महत्तमतेमुळे
  $\lnot\Box\lnot !A \in \Gamma$.
  आता $\Box\lnot!A \in \Gamma$
  असेल, तर गृहीतकानुसार
  $\lnot!A \in \Delta$; पण
  $!A \in \Delta$ असल्याने हे
  $\Delta$ च्या सुसंगततेच्या
  विरुद्ध आहे. म्हणून आवश्यकतेनुसार
  $\Box\lnot!A \notin \Gamma$.

  उजवीकडून डावीकडे जाणारी दिशा: $\Diamond\Delta
  \subseteq \Gamma$ असे समजा.
  प्रतिपरिवर्तन वापरू: $\Box!A
  \notin \Gamma$ दाखवण्यासाठी
  $!A \notin \Delta$ असे समजा.
  $!A \notin \Delta$ असेल,
  तर महत्तमतेमुळे
  $\lnot!A \in \Delta$;
  म्हणून गृहीतकानुसार
  $\Diamond\lnot!A \in \Gamma$.
  पण सामान्य मोडल तर्कशास्त्रात
  $\Diamond\lnot!A$ हे
  $\lnot\Box !A$ शी समतुल्य आहे;
  आणि दुसरे सूत्र~$\Gamma$ मध्ये
  असेल, तर सुसंगततेमुळे
  $\Box!A \notin\Gamma$,
  हेच दाखवायचे होते.
\end{proof}

\iftag{prvBox}{%
  \iftag{prvDiamond}{%
% Box and Diamond are both primitive; prove
% prop:diamond using lem:box-iff-diamond

\begin{prop}\ollabel{prop:diamond}
  $\Gamma$ हा संपूर्ण $\Sigma$-सुसंगत
  असेल, तर $\Diamond !A \in \Gamma$
  असेल तेव्हा आणि तेव्हाच
  $\Diamond\Delta \subseteq
  \Gamma$ असा एखादा संपूर्ण
  $\Sigma$-सुसंगत $\Delta$
  असेल की $!A \in \Delta$.
\end{prop}

\begin{proof}
$\Gamma$ हा संपूर्ण $\Sigma$-सुसंगत
आहे असे समजा. \Dual{} आणि निगामी बंदतेनुसार
$\Diamond!A \in \Gamma$ असेल तेव्हा
आणि तेव्हाच $\lnot\Box\lnot !A \in
\Gamma$. $\Gamma$ संपूर्ण
$\Sigma$-सुसंगत असल्याने
\olref[ccs]{prop:ccs-properties}\olref[ccs]{prop:ccs-lnot}
नुसार $\lnot\Box\lnot !A \in \Gamma$
असेल तेव्हा आणि तेव्हाच
$\Box\lnot !A \notin \Gamma$.
\olref{prop:box} नुसार
$\Box\lnot !A \notin \Gamma$
असेल तेव्हा आणि तेव्हाच
$\Box^{-1}\Gamma \subseteq \Delta$
आणि $\lnot!A \notin \Delta$
असा एखादा संपूर्ण $\Sigma$-सुसंगत
$\Delta$ असतो. असा कोणताही~$\Delta$
घ्या. \olref{lem:box-iff-diamond}
नुसार $\Box^{-1}\Gamma \subseteq \Delta$
असेल तेव्हा आणि तेव्हाच
$\Diamond\Delta \subseteq \Gamma$.
तसेच
\olref[ccs]{prop:ccs-properties}\olref[ccs]{prop:ccs-lnot}
नुसार $\lnot !A \notin \Delta$
असेल तेव्हा आणि तेव्हाच
$!A \in \Delta$. म्हणून
$\Diamond!A \in \Gamma$
असेल तेव्हा आणि तेव्हाच
$\Diamond\Delta \subseteq \Gamma$
आणि $!A \in \Delta$ असा एखादा
संपूर्ण $\Sigma$-सुसंगत $\Delta$
असतो.
\end{proof}

}{}}{}

\begin{prob}
  $\Gamma$ हा संपूर्ण $\Sigma$-सुसंगत
  असेल, तर पुढील विधान दाखवा:
  \iftag{prvBox}{$\Diamond !A \in \Gamma$
    असेल तेव्हा आणि तेव्हाच
    $\Box^{-1}\Gamma \subseteq \Delta$
    आणि $!A \in \Delta$ असा संपूर्ण
    $\Sigma$-सुसंगत $\Delta$ असतो.}
    {$\Box !A \in
    \Gamma$ असेल तेव्हा आणि तेव्हाच
    $\Box^{-1} \Gamma \subseteq \Delta$
    असलेल्या प्रत्येक संपूर्ण
    $\Sigma$-सुसंगत $\Delta$ साठी
    $!A \in \Delta$ असते.}
  \emph{हे
    \olref[nml][com][mod]{lem:box-iff-diamond}
    न वापरता करा}.
\end{prob}

\end{document}
